\documentclass[11pt]{article}
\usepackage[margin=1in]{geometry}
\usepackage{amsmath,amssymb,amsthm}
\usepackage[T1]{fontenc}
\usepackage{lmodern}
\usepackage{hyperref}
\hypersetup{colorlinks=true,linkcolor=blue,citecolor=blue,urlcolor=blue,
 pdftitle={An exponential lower bound for integration under uniform derivative bounds},
 pdfauthor={Kaiwen Jin}}
\newtheorem{theorem}{Theorem}
\newtheorem{lemma}{Lemma}
\newcommand{\R}{\mathbb R}
\newcommand{\N}{\mathbb N}
\newcommand{\tr}{\operatorname{tr}}
\newcommand{\Id}{\operatorname{Id}}
\newcommand{\op}{\mathrm{op}}
\title{An exponential lower bound for integration\protect\\under uniform derivative bounds}
\author{Kaiwen Jin\\\small\href{mailto:yui.ui.siki@gmail.com}{yui.ui.siki@gmail.com}}
\date{October 13, 2026}

\begin{document}
\maketitle
\begin{abstract}
We consider integration on the unit cube for smooth functions whose every
partial derivative, including the function itself, is bounded by one.
For every fixed absolute error tolerance $0<\varepsilon<1$, we prove that
deterministic integration requires at least
$(1-\varepsilon)(\tfrac12+\sin\tfrac12)^{-d}$ function evaluations in
dimension $d$. The bound includes adaptive sampling and nonlinear
reconstruction. The proof uses quadratic forms of tensor products of unit
vectors: a projection makes the function vanish at the samples, while its
integral stays large and all derivative bounds are preserved.
\end{abstract}

\section{The integration problem}

Uniform control of derivatives gives strong local regularity, but the cost
of integrating such functions in high dimension is a separate question.
For each integer $d\ge1$, let
\begin{equation}\label{eq:class}
 F_d=\left\{f\in C^\infty([0,1]^d;\R):
       \|D^\alpha f\|_\infty\le1\text{ for every }\alpha\in\N_0^d\right\},
 \qquad I_d(f)=\int_{[0,1]^d}f(x)\,dx.
\end{equation}
Here $D^\alpha$ denotes the mixed partial derivative with multi-index
$\alpha$, and the supremum is taken on the closed cube. In particular,
$\alpha=0$ is included. Let $n(\varepsilon,d)$ be the smallest number of
exact function evaluations needed by a deterministic algorithm with
worst-case absolute error at most $\varepsilon$ on $F_d$. The sites may be
chosen adaptively and the output may depend nonlinearly on the sampled values.

Wojtaszczyk~\cite[Theorem 2.1]{Wojtaszczyk} proved that, for any fixed number
of samples, the minimal integration error on this class tends to one as
$d\to\infty$. The exponential dependence on dimension remained a distinct
question. Krieg and Ullrich~\cite[Section 8.1, equation (8.5) and the
discussion after Proposition 8.3]{KriegUllrich} explicitly distinguish it
from the known curse for uniform approximation on the same class.
Novak and Pillichshammer~\cite[Section 5.2]{NovakPillichshammer} give related
integration lower bounds for finite-summability derivative norms and
positive quadrature formulas; their discussion leaves the supremum
derivative norm unresolved.

Positive-matrix and squared-kernel arguments give exponential integration
lower bounds in tensor Hilbert spaces~\cite{Vybiral,HinrichsEtAl}.
Those Hilbert unit balls need not satisfy \eqref{eq:class}; the
trigonometric example in~\cite[Section 3.1]{HinrichsEtAl} has a sup-norm
embedding constant $2^{d/2}$. Here we use the operator norm of a projection
to bound a quadratic function directly. Choosing the vector frequency to
be $1/2$ preserves every derivative bound, and a tensor moment calculation
retains the absolute integration scale.

\begin{theorem}\label{thm:main}
Set $\lambda=\tfrac12+\sin(\tfrac12)$, with angles in radians. For every
$d\ge1$ and $n\in\N_0$, any deterministic algorithm using at most $n$
exact function evaluations has worst-case absolute integration error on
$F_d$ at least
\begin{equation}\label{eq:error}
 \max\{0,1-n\lambda^d\}.
\end{equation}
Consequently, for every $0<\varepsilon<1$ and $d\ge1$,
\begin{equation}\label{eq:complexity}
 n(\varepsilon,d)\ge (1-\varepsilon)\lambda^{-d}.
\end{equation}
This includes adaptive algorithms and nonlinear reconstruction, without
any restriction on quadrature weights.
\end{theorem}

Since $\tfrac12<\lambda<1$, equation~\eqref{eq:complexity} gives the curse
of dimensionality: one may take $\varepsilon_0=1/2$, $c=1/2$ and
$\rho=\lambda^{-1}>1$, in every dimension.

\section{Proof}

We construct a bounded function that vanishes at any prescribed sample
sites. Its values are squared distances from a subspace in a finite
Euclidean space. The vector defining this distance has derivatives of
exact norm $2^{-k}$ at order $k$, which offsets the factors in the product rule.

For $t\in\R$ and $x\in\R^d$, define
\begin{equation}\label{eq:vectors}
 u(t)=\begin{pmatrix}\cos((t-1/2)/2)\\\sin((t-1/2)/2)\end{pmatrix},
 \qquad U_d(x)=\bigotimes_{j=1}^d u(x_j)\in\R^{2^d}.
\end{equation}
All vector norms below are Euclidean. The tensor product in
\eqref{eq:vectors} has one coordinate for each choice of a sine or cosine
component in every factor.

\begin{lemma}\label{lem:derivatives}
If $A$ is a real symmetric $2^d\times2^d$ matrix with
$\|A\|_\op\le1$, then
$f_A(x)=U_d(x)^T A U_d(x)$ belongs to $F_d$.
\end{lemma}

\begin{proof}
Let $J=\left(\begin{smallmatrix}0&-1\\1&0\end{smallmatrix}\right)$.
Then $u^{(k)}(t)=2^{-k}J^k u(t)$ and $\|u(t)\|_2=1$, so
$\|u^{(k)}(t)\|_2=2^{-k}$ for all $k\ge0$. Each coordinate derivative
acts on its own tensor factor. The squared norm of a tensor product is the
product of the squared factor norms; therefore
\begin{equation}\label{eq:vector-derivatives}
 \|D^\beta U_d(x)\|_2=2^{-|\beta|},\qquad
 |\beta|=\beta_1+\cdots+\beta_d.
\end{equation}
For every $\alpha\in\N_0^d$, the product rule gives
\[
 D^\alpha f_A(x)=\sum_{\beta\le\alpha}
 {\alpha\choose\beta}(D^\beta U_d(x))^T A D^{\alpha-\beta}U_d(x),
 \qquad {\alpha\choose\beta}=\prod_{j=1}^d{\alpha_j\choose\beta_j},
\]
where $\beta\le\alpha$ means coordinatewise comparison. By
Cauchy--Schwarz and \eqref{eq:vector-derivatives},
\[
 |D^\alpha f_A(x)|\le
 2^{-|\alpha|}\sum_{\beta\le\alpha}{\alpha\choose\beta}
 =2^{-|\alpha|}\prod_{j=1}^d2^{\alpha_j}=1.
\]
The finite sums and products defining $f_A$ are smooth on all of $\R^d$.
The estimate includes order zero and every repeated or mixed derivative.
\end{proof}

\begin{lemma}\label{lem:fooler}
For any $m\in\N_0$ and sites $x^1,\ldots,x^m\in[0,1]^d$, there is a
function $f\in F_d$ such that $0\le f\le1$, $f(x^j)=0$ for every $j$, and
\begin{equation}\label{eq:integral-fooler}
 I_d(f)\ge1-m\lambda^d.
\end{equation}
\end{lemma}

\begin{proof}
Let $P$ be the orthogonal projection onto
$V=\operatorname{span}\{U_d(x^1),\ldots,U_d(x^m)\}$ and let $Q=\Id-P$.
For $m=0$, take $P=0$. The projection $Q$ is a symmetric contraction, so
Lemma~\ref{lem:derivatives} applies to
\[
 f(x)=U_d(x)^TQ U_d(x)=\|Q U_d(x)\|_2^2.
\]
Thus $0\le f\le1$ and $f(x^j)=0$.

To estimate its integral, compute the one-dimensional moment matrix
\begin{equation}\label{eq:moment}
 M=\int_0^1u(t)u(t)^T\,dt
 =\begin{pmatrix}\lambda&0\\0&1-\lambda\end{pmatrix}.
\end{equation}
The off-diagonal entry vanishes by oddness about $1/2$, and
\[
 \int_0^1\cos^2((t-1/2)/2)\,dt
 =\frac12+\frac12\int_0^1\cos(t-1/2)\,dt
 =\frac12+\sin(1/2).
\]
The sine-square entry is its complement. Since $0<\sin(1/2)<1/2$,
$\tfrac12<\lambda<1$. Product integration of the continuous entries yields
\[
 M_d=\int_{[0,1]^d}U_d(x)U_d(x)^T\,dx=M^{\otimes d}.
\]
The matrix $M_d$ is diagonal, has trace one, and every diagonal entry is
at most $\lambda^d$. Write $r=\dim V\le m$ and choose an orthonormal
basis $v_1,\ldots,v_r$ of $V$. Then $P=\sum_{i=1}^r v_i v_i^T$, and
\[
 \tr(PM_d)=\sum_{i=1}^r v_i^T M_d v_i\le r\lambda^d\le m\lambda^d.
\]
It follows that
$I_d(f)=\tr(QM_d)=1-\tr(PM_d)\ge1-m\lambda^d$.
Repeated or dependent sample sites simply reduce $r$.
\end{proof}

\begin{proof}[Proof of Theorem~\ref{thm:main}]
Run the algorithm on the zero function and record its $m\le n$ sample
sites, including its stopping decision. Apply Lemma~\ref{lem:fooler} to
these sites. Both $f$ and $-f$ belong to $F_d$ and return zero at every
recorded site. By induction over the sampling steps, the algorithm follows
the same path on either input and produces the same output $a$.
For $b=I_d(f)\ge0$, the triangle inequality gives
\[
 \max\{|b-a|,|-b-a|\}\ge b\ge1-m\lambda^d\ge1-n\lambda^d.
\]
This proves \eqref{eq:error}, since an error is nonnegative. If an
algorithm has worst-case error at most $\varepsilon$, then
$1-n\lambda^d\le\varepsilon$, which is \eqref{eq:complexity}.
\end{proof}

\paragraph{AI use.}
OpenAI Codex assisted with the literature search, construction and checking
of the argument, and preparation of the manuscript. Independent
mathematical review is requested.

\begin{thebibliography}{9}
\bibitem{KriegUllrich}
D. Krieg and M. Ullrich,
\emph{Approximation of functions: Optimal sampling and complexity},
Acta Numerica \textbf{35} (2026), 273--457.
\href{https://doi.org/10.1017/S0962492925100287}{doi:10.1017/S0962492925100287}.
Preprint version checked:
\href{https://arxiv.org/abs/2602.02066v2}{arXiv:2602.02066v2}.

\bibitem{Wojtaszczyk}
J. O. Wojtaszczyk,
\emph{Multivariate integration in $C^\infty([0,1]^d)$ is not strongly tractable},
Journal of Complexity \textbf{19} (2003), 638--643.
\href{https://arxiv.org/abs/0803.0441v1}{arXiv:0803.0441v1}.

\bibitem{NovakPillichshammer}
E. Novak and F. Pillichshammer,
\emph{Intractability results for integration in tensor product spaces},
preprint (2024),
\href{https://arxiv.org/abs/2404.17163v1}{arXiv:2404.17163v1}.

\bibitem{Vybiral}
J. Vyb\'iral,
\emph{A variant of Schur's product theorem and its applications},
Advances in Mathematics \textbf{368} (2020), 107140.
\href{https://doi.org/10.1016/j.aim.2020.107140}{doi:10.1016/j.aim.2020.107140}.

\bibitem{HinrichsEtAl}
A. Hinrichs, D. Krieg, E. Novak and J. Vyb\'iral,
\emph{Lower bounds for the error of quadrature formulas for Hilbert spaces},
Journal of Complexity \textbf{65} (2021), 101544.
\href{https://arxiv.org/abs/2004.00274v2}{arXiv:2004.00274v2}.
\end{thebibliography}
\end{document}
